Lithium-ion batteries are used extensively in transportation, portable
electronics, aerospace systems, stationary energy storage, and numerous
other applications requiring high energy density and rechargeable
electrochemical storage.

Despite their advantages, lithium-ion cells can undergo thermal runaway
when subjected to sufficiently severe thermal, electrical, mechanical,
or internal abuse conditions. Thermal runaway is characterized by a
self-accelerating sequence of exothermic reactions that can produce rapid
temperature increase, gas generation, internal pressurization, venting,
fire, and, under some conditions, propagation to neighbouring cells.

The physical processes involved span several strongly coupled domains,
including

\begin{itemize}
    \item electrochemical state evolution;
    \item decomposition reaction kinetics;
    \item heat generation and heat transfer;
    \item gas production;
    \item thermochemical gas-phase reactions;
    \item internal pressure evolution;
    \item venting and structural failure;
    \item external combustion;
    \item cell-to-cell thermal propagation.
\end{itemize}

A single mathematical model rarely represents all of these phenomena
with the same level of fidelity. Thermal runaway simulations therefore
often rely on reduced-order models calibrated against calorimetric,
abuse-test, or literature data.

LiionTR is being developed as a modular scientific framework for
constructing and solving such models while preserving a clear separation
between individual physical mechanisms.


\subsection{Project objectives}

The primary objective of LiionTR is to provide a reusable computational
framework for lithium-ion battery thermal runaway research.

The framework is designed to support:

\begin{itemize}
    \item multiple thermal runaway reaction mechanisms;
    \item reusable kinetic and reaction-progress laws;
    \item configurable battery-cell geometries and materials;
    \item thermal energy balances;
    \item empirical and future thermochemical gas-generation models;
    \item internal gas-pressure calculations;
    \item compressible venting;
    \item numerical event handling;
    \item model verification and validation;
    \item future cell-to-cell propagation models.
\end{itemize}


\subsection{Design philosophy}

A central design principle is the separation of physical models from
numerical integration.

For example, an Arrhenius model defines a temperature-dependent kinetic
rate but does not know how the corresponding reaction state is
integrated.

Similarly, a pressure model evaluates pressure from thermodynamic state
variables but does not advance those variables in time.

The numerical solver acts as the orchestration layer that assembles these
individual components into the complete governing system.

This separation is intended to improve

\begin{itemize}
    \item physical transparency;
    \item software maintainability;
    \item testability;
    \item reuse of scientific models;
    \item comparison between alternative formulations.
\end{itemize}


\subsection{Reduced-order modeling}

The current LiionTR formulation is primarily zero-dimensional.

A battery cell is represented using a lumped temperature state and a set
of scalar reaction-progress variables.

The thermal runaway process is described conceptually by

\begin{equation}
    \text{temperature}
    \rightarrow
    \text{reaction rates}
    \rightarrow
    \text{heat generation}
    \rightarrow
    \text{temperature},
\end{equation}

forming the fundamental positive feedback loop responsible for thermal
runaway.

Additional coupling can be introduced through

\begin{equation}
    \text{reaction progress}
    \rightarrow
    \text{gas generation}
    \rightarrow
    \text{pressure}
    \rightarrow
    \text{venting}.
\end{equation}

These reduced-order models are computationally inexpensive compared with
fully spatial electrochemical, thermal, mechanical, or computational
fluid dynamics models. They are therefore useful for parameter studies,
model development, sensitivity analysis, and safety-oriented system
simulations.


\subsection{Current implementation}

The present framework includes:

\begin{itemize}
    \item generic cell, geometry, and material representations;
    \item Arrhenius and threshold kinetic models;
    \item multiple reaction-progress formulations;
    \item context-dependent reaction coupling;
    \item standard and multi-channel reactions;
    \item reaction networks;
    \item a lumped thermal model;
    \item empirical reaction-based gas generation;
    \item species-resolved gas inventories;
    \item an ideal-gas pressure model;
    \item choked and unchoked compressible vent flow;
    \item irreversible vent-opening events;
    \item stiff numerical integration.
\end{itemize}

The framework also includes an implementation of the reduced-order
thermal runaway reaction model reported by
\textcite{hu2020inhibition}.


\subsection{Relationship to external scientific software}

LiionTR is not intended to reproduce capabilities that are already
provided by specialized scientific software.

Instead, external scientific packages may eventually be incorporated as
backends.

Two particularly relevant examples are Cantera and PyBaMM.

Cantera provides thermodynamic, chemical-kinetic, and transport
capabilities and is a candidate backend for equilibrium gas composition
and temperature-dependent mixture properties
\parencite{cantera2025}.

PyBaMM provides a framework for physics-based battery modeling and may
supply electrochemical states ranging from reduced-order models to
porous-electrode formulations
\parencite{sulzer2021pybamm}.

The intended long-term relationship is therefore

\begin{equation}
    \text{PyBaMM}
    \rightarrow
    \text{electrochemical state},
\end{equation}

\begin{equation}
    \text{Cantera}
    \rightarrow
    \text{thermochemical state},
\end{equation}

and

\begin{equation}
    \text{LiionTR}
    \rightarrow
    \text{thermal runaway and safety coupling}.
\end{equation}


\subsection{Scope of this report}

This report documents the mathematical formulation and current software
architecture of LiionTR.

The following sections describe:

\begin{enumerate}
    \item the general physical and mathematical framework;
    \item the lumped thermal formulation;
    \item reaction kinetics and progress models;
    \item reaction-network construction;
    \item implementation of the Hu et al.\ thermal runaway model;
    \item gas generation;
    \item pressure and venting;
    \item numerical integration;
    \item verification strategy;
    \item planned future development.
\end{enumerate}

The report describes the current state of the framework and therefore
distinguishes between implemented capabilities and planned future
extensions.